\section{A tétel lényege és vizsgaszintű vázlata}

\begin{summarybox}
A 23. tétel az adatkapcsolati réteg közeghozzáférési alrétegéről szól. A központi kérdés az, hogy több állomás hogyan osztozhat egy közös üzenetszórásos csatornán. Ebből következnek a statikus és dinamikus csatornamegosztási módszerek, a versengő és determinisztikus eljárások, majd két fontos gyakorlati példa: az IEEE 802.3 Ethernet CSMA/CD eljárása és az IEEE 802.11 WLAN CSMA/CA alapú működése. Vizsgán különösen fontos az Ethernet-keret, a MAC-cím, a half duplex és full duplex Ethernet különbsége, valamint a WLAN-nál az AP, BSS, ESS, DCF, PCF, RTS/CTS és NAV szerepe.
\end{summarybox}

Vizsgán célszerű a következő sorrendet követni:

\begin{enumerate}
  \item adatkapcsolati réteg két alrétege: LLC és MAC;
  \item a MAC alréteg feladata: közös üzenetszórásos csatorna megosztása;
  \item csatornamegosztási módszerek: statikus és dinamikus;
  \item dinamikus módszerek: versengő/ütközéses és determinisztikus/ütközésmentes;
  \item ALOHA, réselt ALOHA, CSMA, CSMA/CD és CSMA/CA rövid összehasonlítása;
  \item IEEE 802 szabványcsalád és az IEEE 802.3 Ethernet helye;
  \item Ethernet CSMA/CD működése, ütközésérzékelés, jam jel, exponenciális visszatartás;
  \item Ethernet/IEEE 802.3 keretformátum, minimális és maximális keretméret;
  \item MAC-címek: unicast, multicast, broadcast, globális és lokális címzés;
  \item támogatott Ethernet közegek és sebességek, half duplex és full duplex működés;
  \item IEEE 802.11 WLAN felépítése: BSS, AP, DS, ESS, roaming;
  \item CSMA/CA, DCF, pozitív nyugta, interframe space-ek, backoff, RTS/CTS és NAV;
  \item PCF mint központosított, versengésmentes kiegészítő üzemmód;
  \item gyakori vizsgahibák és rövid összefoglaló válasz.
\end{enumerate}

A legfontosabb vizsgamondat: a MAC alréteg üzenetszórásos vagy több hozzáférésű közeg esetén azt szabályozza, hogy az állomások mikor és milyen feltételekkel használhatják a közös csatornát; az Ethernet klasszikusan CSMA/CD-t, a WLAN pedig a rádiós közeg sajátosságai miatt CSMA/CA-t alkalmaz.

\section{Az adatkapcsolati réteg és a MAC alréteg}

\subsection{LLC és MAC szerepe}

Az ISO--OSI modellben az \textbf{adatkapcsolati réteg} a fizikai réteg által biztosított nyers bitfolyam fölött működik. Feladata, hogy a szomszédos csomópontok között kereteket vigyen át, felismerje a kerethatárokat, ellenőrizze a hibákat, szükség esetén adatfolyam-vezérlést végezzen, és a hálózati réteg felé szabályos szolgáltatást nyújtson.

Az IEEE 802 szemléletben az adatkapcsolati réteg két alrétegre bontható:

\begin{itemize}
  \item \textbf{LLC -- Logical Link Control}: a felsőbb rétegek felé egységes logikai kapcsolatvezérlési szolgáltatást ad;
  \item \textbf{MAC -- Media Access Control}: a fizikai közeghez való hozzáférést szabályozza.
\end{itemize}

Pont--pont kapcsolat esetén a közeg megosztásának kérdése általában egyszerű, mert csak két végpont használja a vonalat. Üzenetszórásos vagy többszörös hozzáférésű közeg esetén viszont több állomás ugyanazon a csatornán próbálhat adni, ezért külön közeghozzáférési szabályokra van szükség.

\subsection{A MAC alréteg feladata}

A \textbf{MAC alréteg} feladata egy közös üzenetszórásos csatorna megosztása több egymással versengő állomás között. Olyan szabályrendszert ad, amely eldönti, hogy adott időpillanatban melyik állomás adhat, hogyan lehet ütközést felismerni vagy elkerülni, és hogyan kell hiba vagy ütközés után újrapróbálkozni.

A MAC alréteg célja az, hogy a közös közeg felett a felsőbb, LLC alréteg számára pont--pont jellegű, keretekkel dolgozó szolgáltatást biztosítson. Ennek sikere a forgalom jellegétől is függ. Folytonos, egyenletes forgalomnál más módszer kedvező, mint löketszerű, véletlenszerűen jelentkező keretküldési igényeknél.


\begin{figure}[htbp]
\centering
\resizebox{0.88\textwidth}{!}{%
\begin{tikzpicture}[font=\scriptsize, node distance=5mm]
  \tikzset{layer/.style={draw, rounded corners, minimum width=45mm, minimum height=8mm, align=center}}
  \node[layer, fill=black!5] (net) {3. Hálózati réteg\\csomagok, útválasztás};
  \node[layer, fill=lightaccent, below=of net] (llc) {LLC alréteg\\logikai kapcsolatvezérlés};
  \node[layer, fill=warnbg, below=of llc] (mac) {MAC alréteg\\közeghozzáférés vezérlése};
  \node[layer, fill=black!5, below=of mac] (phy) {1. Fizikai réteg\\jelek, bitfolyam, közeg};

  \node[layer, fill=lightaccent, right=22mm of mac, minimum width=48mm] (shared) {Üzenetszórásos közeg\\több állomás versenghet};
  \draw[arrow] (net) -- (llc);
  \draw[arrow] (llc) -- (mac);
  \draw[arrow] (mac) -- (phy);
  \draw[arrow] (mac.east) -- (shared.west) node[midway, above]{csatornamegosztás};

  \node[draw, rounded corners, fill=black!5, below=7mm of shared, align=center, minimum width=48mm] (goal)
    {MAC célja:\\pont--pont jellegű szolgáltatás\\közös közeg felett};
  \draw[arrow] (shared) -- (goal);
\end{tikzpicture}%
}
\caption{Az adatkapcsolati réteg LLC és MAC alrétege}
\label{fig:zv23_mac_alretegek}
\end{figure}


\section{Csatornamegosztási módszerek}

\subsection{Statikus és dinamikus csatornamegosztás}

A közös csatorna megosztási módjai két nagy csoportba sorolhatók.

\textbf{Statikus csatornamegosztásnál} a csatorna kapacitását előre, rögzített módon osztják fel a felhasználók között. Ilyen például az FDM és a TDM. FDM esetén minden felhasználó külön frekvenciasávot kap, TDM esetén pedig minden felhasználóhoz egy vagy több időrés tartozik. A módszer jól működik egyenletes, tartós forgalomnál, viszont löketszerű forgalomnál pazarló, mert a kiosztott rész akkor is üresen marad, ha az adott állomásnak éppen nincs küldenivalója.

\textbf{Dinamikus csatornamegosztásnál} a csatorna azok között oszlik meg, akiknek éppen küldenivalójuk van. A döntés lehet központosított vagy elosztott, a hozzáférés pedig lehet versengő vagy determinisztikus.

\begin{center}

\begin{tabularx}{\textwidth}{L{30mm}L{30mm}X}
\toprule
\textbf{Módszer} & \textbf{Példa} & \textbf{Jellemzés} \\
\midrule
Statikus & FDM, TDM & A csatorna kapacitását előre kiosztják. Egyenletes, tartós forgalomnál jó, löketszerű forgalomnál pazarló lehet. \\
Dinamikus centralizált & polling, PCF & Egy kijelölt állomás vagy koordinátor dönt a hozzáférésről. Jobban szabályozható, de központi elemhez kötött. \\
Dinamikus elosztott versengő & ALOHA, CSMA, CSMA/CD, CSMA/CA & Az állomások versengenek a csatornáért. Kis terhelésnél gyors, nagy terhelésnél ütközések és visszatartások miatt romolhat. \\
Dinamikus determinisztikus & vezérjel, foglalás & Szabályozott, ütközésmentes hozzáférést ad. Korlátozható a várakozás, de a vezérlési többlet nagyobb. \\
\bottomrule
\end{tabularx}

\end{center}


\begin{figure}[htbp]
\centering
\resizebox{0.95\textwidth}{!}{%
\begin{tikzpicture}[font=\scriptsize, node distance=8mm]
  \tikzset{box/.style={draw, rounded corners, align=center, minimum width=35mm, minimum height=8mm}}
  \node[box, fill=accent!12] (root) {Csatornamegosztás};
  \node[box, fill=lightaccent, below left=13mm and 22mm of root] (stat) {Statikus\\előre kiosztott rész};
  \node[box, fill=lightaccent, below right=13mm and 22mm of root] (dyn) {Dinamikus\\aktuális igény szerint};
  \node[box, fill=black!5, below left=11mm and 6mm of stat] (fdm) {FDM\\frekvenciaosztás};
  \node[box, fill=black!5, below right=11mm and 6mm of stat] (tdm) {TDM\\időosztás};
  \node[box, fill=warnbg, below left=11mm and 6mm of dyn] (cont) {Versengő\\ütközéses};
  \node[box, fill=warnbg, below right=11mm and 6mm of dyn] (det) {Determinisztikus\\ütközésmentes};
  \node[box, fill=black!5, below=10mm of cont] (aloha) {ALOHA, CSMA,\\CSMA/CD, CSMA/CA};
  \node[box, fill=black!5, below=10mm of det] (token) {vezérjel, polling,\\foglalás};

  \draw[arrow] (root) -- (stat);
  \draw[arrow] (root) -- (dyn);
  \draw[arrow] (stat) -- (fdm);
  \draw[arrow] (stat) -- (tdm);
  \draw[arrow] (dyn) -- (cont);
  \draw[arrow] (dyn) -- (det);
  \draw[arrow] (cont) -- (aloha);
  \draw[arrow] (det) -- (token);
\end{tikzpicture}%
}
\caption{Csatornamegosztási módszerek fő csoportjai}
\label{fig:zv23_csatornamegosztas}
\end{figure}


\subsection{Versengő és determinisztikus dinamikus módszerek}

\textbf{Versengő}, más néven ütközéses módszernél az állomások nem kapnak előre garantált küldési időpontot. Amikor küldenivalójuk van, megpróbálják megszerezni a csatornát. Ha több állomás egyszerre ad, ütközés történhet. Ide tartozik az ALOHA, a réselt ALOHA, a CSMA, a CSMA/CD és a CSMA/CA.

A versengő módszerek előnye, hogy kis terhelésnél alacsony késleltetést adhatnak, mert az állomás gyorsan adhat, ha a csatorna szabad. Hátrányuk, hogy nagy terhelésnél sok lehet az ütközés vagy az újrapróbálkozás, ezért a csatorna kihasználtsága romolhat.

\textbf{Determinisztikus}, más néven ütközésmentes módszernél a hozzáférés szabályozottabb. Például vezérjel, lekérdezés vagy foglalási mechanizmus döntheti el, ki adhat. Előnye, hogy korlátozható a maximális várakozási idő és elkerülhető az ütközés. Hátránya, hogy bonyolultabb, és kis forgalomnál a vezérlési többlet aránylag nagy lehet.

\subsection{ALOHA, CSMA, CSMA/CD és CSMA/CA}

Az \textbf{ALOHA} alapgondolata egyszerű: az állomás akkor ad, amikor elkészült a kerete. Ha nincs nyugta, akkor feltételezi, hogy ütközés vagy hiba történt, véletlen ideig vár, majd újraküldi. A réselt ALOHA javítás: az adás csak időrések elején kezdhető, ezért kevesebb az ütközési lehetőség.

A \textbf{CSMA} -- Carrier Sense Multiple Access -- csatornafigyelő többszörös hozzáférés. Az állomás adás előtt belehallgat a csatornába. Ha szabad, adni kezd; ha foglalt, vár. Több változata van:

\begin{itemize}
  \item \textbf{1-perzisztens CSMA}: ha a csatorna felszabadul, az állomás azonnal adni kezd;
  \item \textbf{nemperzisztens CSMA}: ha a csatorna foglalt, véletlen ideig vár, majd újra figyel;
  \item \textbf{p-perzisztens CSMA}: réselt csatornán tétlen közeg esetén p valószínűséggel ad, különben vár a következő résig.
\end{itemize}

A \textbf{CSMA/CD} a CSMA ütközésérzékeléssel kiegészített változata. Az állomás adás közben is figyeli a közeget; ha ütközést érzékel, megszakítja az adást, jam jelet küld, majd véletlen visszatartás után újrapróbálkozik. Ez volt a klasszikus megosztott közegű Ethernet alapja.

A \textbf{CSMA/CA} a CSMA ütközéselkerüléssel kiegészített változata. Rádiós hálózatban az állomás adás közben nem tud megbízhatóan ütközést érzékelni, ezért inkább megpróbálja csökkenteni az ütközések valószínűségét. Ilyenkor a közeg figyelése, az interframe space idők, a véletlen backoff, a pozitív ACK és szükség esetén az RTS/CTS mechanizmus együtt valósítja meg a hozzáférést.

\begin{center}

\begin{tabularx}{\textwidth}{L{24mm}L{34mm}X}
\toprule
\textbf{Módszer} & \textbf{Alapelv} & \textbf{Megjegyzés} \\
\midrule
ALOHA & Az állomás azonnal ad, ha kerete van. & Egyszerű, de sok az ütközés; kis terhelésű, sok állomásos környezetben értelmezhető. \\
Réselt ALOHA & Az adás csak időrések elején kezdhető. & Csökkenti az ütközési lehetőséget, de szinkronizációt igényel. \\
CSMA & Adás előtt csatornafigyelés. & Jobb kihasználtság, de a terjedési késleltetés miatt továbbra is lehet ütközés. \\
CSMA/CD & CSMA + CD. & Klasszikus megosztott közegű Ethernet; ütközéskor jam jel és backoff következik. \\
CSMA/CA & CSMA ütközéselkerüléssel. & WLAN-oknál használatos; DIFS, SIFS, ACK, backoff, NAV és opcionális RTS/CTS segíti. \\
\bottomrule
\end{tabularx}

\end{center}

\section{IEEE 802 és az Ethernet helye}

\subsection{IEEE 802 szabványcsalád}

Az \textbf{IEEE 802} a helyi és városi hálózatokra vonatkozó szabványcsalád. Az ISO is elfogadta ISO 8802 néven. A szabványcsalád részei közül a 802.2 az LLC alréteget, a különböző 802.x részek pedig tipikusan MAC- és fizikai rétegbeli technológiákat írnak le.

Fontosabb példák:

\begin{itemize}
  \item \textbf{802.3}: Ethernet, CSMA/CD alapú történeti megosztott közegű LAN;
  \item \textbf{802.4}: Token Bus;
  \item \textbf{802.5}: Token Ring;
  \item \textbf{802.11}: vezeték nélküli LAN, WLAN;
  \item \textbf{802.16}: vezeték nélküli MAN.
\end{itemize}

Az Ethernet eredetileg a Xerox, DEC és Intel fejlesztése volt. Az IEEE 802.3 az Ethernet szabványosított változata. Az Ethernet II és az IEEE 802.3 között van keretformátumbeli különbség, de a klasszikus közeghozzáférési alapelv közös: a megosztott közegű Ethernet CSMA/CD-t használ.

\subsection{Ethernet fizikai közegek és jelölések}

Az Ethernet jelölések általános alakja például \textbf{10BASE5}. A szám az átviteli sebességet jelöli Mbit/s-ban, a BASE az alapsávi átvitelt, az utolsó rész pedig a közeget vagy a maximális szegmenshosszt. A klasszikus példák:

\begin{itemize}
  \item \textbf{10BASE5}: vastag koaxiális kábel, 500 m szegmenshossz;
  \item \textbf{10BASE2}: vékony koaxiális kábel, körülbelül 185 m szegmenshossz;
  \item \textbf{10BASE-T}: csavart érpár, pont--pont kapcsolat, jellemzően 100 m;
  \item \textbf{10BASE-F}: optikai szál;
  \item \textbf{100BASE-TX/FX}: Fast Ethernet;
  \item \textbf{1000BASE-SX/LX/CX/T}: Gigabit Ethernet különböző optikai és réz közegeken;
  \item \textbf{10GBASE-*}: 10 Gigabit Ethernet, kezdetben főleg optikai, később réz alapú változatokkal is.
\end{itemize}

\begin{center}

\begin{tabularx}{\textwidth}{L{30mm}L{24mm}L{32mm}X}
\toprule
\textbf{Változat} & \textbf{Sebesség} & \textbf{Közeg} & \textbf{Jellemző} \\
\midrule
10BASE5 & 10 Mbit/s & vastag koax & Klasszikus sárga kábel, busz jelleg, kb. 500 m szegmenshossz. \\
10BASE2 & 10 Mbit/s & vékony koax & Cheapernet, BNC és T-csatlakozók, kb. 185 m szegmenshossz. \\
10BASE-T & 10 Mbit/s & csavart érpár & Pont--pont kábelezés hubhoz vagy switchhez, jellemzően 100 m. \\
100BASE-TX/FX & 100 Mbit/s & UTP vagy optika & Fast Ethernet, 802.3u, kapcsolt hálózatokban nagyon elterjedt. \\
1000BASE-SX/LX/T & 1 Gbit/s & optika vagy UTP & Gigabit Ethernet; 1000BASE-T tipikusan 100 m Cat5 vagy jobb kábelen. \\
10GBASE-* & 10 Gbit/s & főleg optika, később réz & 10 Gigabit Ethernet; full duplex működés, a keretformátum alapvetően változatlan. \\
\bottomrule
\end{tabularx}

\end{center}

\section{IEEE 802.3 Ethernet és CSMA/CD}

\subsection{CSMA/CD működési elve}

A klasszikus IEEE 802.3 Ethernet \textbf{CSMA/CD} eljárást használ. Az elnevezés elemei:

\begin{itemize}
  \item \textbf{Carrier Sense}: az állomás adás előtt figyeli, hogy tétlen-e a csatorna;
  \item \textbf{Multiple Access}: több állomás használja ugyanazt a közös közeget;
  \item \textbf{Collision Detection}: az állomás adás közben is képes ütközést érzékelni.
\end{itemize}

A küldés menete leegyszerűsítve:

\begin{enumerate}
  \item az állomás összeállítja a keretet;
  \item figyeli a csatornát;
  \item ha a csatorna foglalt, vár, amíg felszabadul;
  \item ha a csatorna tétlen, az interframe gap után adni kezd;
  \item adás közben figyeli, történt-e ütközés;
  \item ütközés esetén megszakítja az adást, jam jelet küld, növeli a kísérletszámot, majd véletlen visszatartási idő után újrapróbálkozik;
  \item ha nincs ütközés és a keret végig lement, a küldés sikeres.
\end{enumerate}


\begin{figure}[htbp]
\centering
\resizebox{0.94\textwidth}{!}{%
\begin{tikzpicture}[font=\scriptsize, node distance=7mm]
  \tikzset{procstep/.style={draw, rounded corners, fill=lightaccent, align=center, minimum width=31mm, minimum height=8mm},
           decision/.style={draw, diamond, aspect=2.2, fill=warnbg, align=center, inner sep=1pt, minimum width=28mm}}
  \node[procstep] (ready) {Keret kész};
  \node[decision, right=10mm of ready] (busy) {Csatorna\\foglalt?};
  \node[procstep, above=8mm of busy] (wait) {Várakozás\\felszabadulásig};
  \node[procstep, right=12mm of busy] (send) {Adás indítása\\IFG után};
  \node[decision, right=12mm of send] (coll) {Ütközés?};
  \node[procstep, below=8mm of coll] (jam) {Jam jel\\kísérletszám nő};
  \node[procstep, left=13mm of jam] (backoff) {Véletlen\\backoff};
  \node[procstep, right=12mm of coll] (ok) {Keret végigment\\sikeres küldés};

  \draw[arrow] (ready) -- (busy);
  \draw[arrow] (busy) -- node[right]{igen} (wait);
  \draw[arrow] (wait.west) -| (ready.north);
  \draw[arrow] (busy) -- node[above]{nem} (send);
  \draw[arrow] (send) -- (coll);
  \draw[arrow] (coll) -- node[above]{nem} (ok);
  \draw[arrow] (coll) -- node[right]{igen} (jam);
  \draw[arrow] (jam) -- (backoff);
  \draw[arrow] (backoff.west) -| (ready.south);
\end{tikzpicture}%
}
\caption{CSMA/CD keretküldés leegyszerűsített folyamata}
\label{fig:zv23_csma_cd_folyamat}
\end{figure}


\subsection{Ütközés, jam jel és slot time}

Ütközés akkor történik, ha két állomás adása a közegben átfed. Mivel a jel terjedési késleltetéssel halad, előfordulhat, hogy egy állomás még szabadnak látja a csatornát, miközben egy távoli állomás már adni kezdett, csak a jel még nem ért oda.

A CSMA/CD lényege, hogy az állomás adás közben is hallgatja a közeget. Ha eltérést érzékel, akkor ütközést állapít meg, megszakítja a hasznos keret adását, és \textbf{jam jelet} küld, hogy a többi állomás is biztosan észrevegye az ütközést.

A minimális keretméretet az indokolja, hogy a legtávolabbi két állomás közötti oda-vissza terjedési idő alatt is érzékelhető legyen az ütközés. Ha a keret túl rövid lenne, az adó már befejezhetné a küldést, mielőtt megtudná, hogy ütközés történt. Klasszikus Ethernetben ezért a minimális keretméret a célcímtől az FCS végéig \textbf{64 bájt}, azaz 512 bit.


\begin{figure}[htbp]
\centering
\resizebox{0.9\textwidth}{!}{%
\begin{tikzpicture}[font=\scriptsize]
  \draw[thick] (0,0) -- (11,0) node[right]{idő};
  \draw[thick] (0,-1.6) -- (11,-1.6) node[right]{idő};
  \node[left] at (0,0) {A állomás};
  \node[left] at (0,-1.6) {B állomás};

  \draw[fill=lightaccent, draw=accent] (0.5,0.15) rectangle (5.6,0.55) node[midway]{A adni kezd};
  \draw[fill=warnbg, draw=orange!80!black] (3.2,-1.45) rectangle (6.2,-1.05) node[midway]{B adni kezd};

  \draw[dashed] (0.5,0.15) -- (2.1,-1.05) node[midway, above, sloped] {$\tau$};
  \draw[dashed] (3.2,-1.05) -- (4.8,0.15) node[midway, below, sloped] {$\tau$};
  \draw[decorate, decoration={brace, amplitude=4pt}] (0.5,0.8) -- (4.8,0.8) node[midway, above=4pt]{legfeljebb $2\tau$ alatt derül ki az ütközés};

  \node[draw, rounded corners, fill=black!5, align=center] at (8.1,-0.55) {Minimális keretidő\\legalább $2\tau$\\klasszikus Ethernetben 512 bit};
\end{tikzpicture}%
}
\caption{Ütközésérzékelés és minimális keretidő CSMA/CD esetén}
\label{fig:zv23_csma_cd_idozites}
\end{figure}


\subsection{Visszatartás és újrapróbálkozás}

Ütközés után az állomások nem próbálkozhatnak azonnal újra, mert akkor jó eséllyel ismét egyszerre indulnának. Ehelyett véletlen visszatartási időt választanak. Az Ethernet \textbf{bináris exponenciális visszatartást} használ: több egymást követő ütközés után a lehetséges várakozási ablak nő, ezért csökken az újabb azonnali ütközés valószínűsége.

Ez a módszer jól illeszkedik a löketszerű LAN-forgalomhoz, de nagy terhelésnél a sok ütközés miatt a hasznos átbocsátóképesség romlik. Ezért a modern Ethernet-telepítésekben a megosztott buszos/hubos half duplex működést gyakorlatilag felváltotta a kapcsolókon alapuló, pont--pont full duplex Ethernet.

\section{Ethernet keretformátum és MAC-címek}

\subsection{IEEE 802.3 és Ethernet II keret}

Az Ethernet keret legfontosabb mezői a preambulum, a keretkezdet jelző, a célcím, a forráscím, a hossz vagy típus mező, az adatmező, szükség esetén a kitöltés, valamint az FCS. A preambulum és az SFD a vevő szinkronizációját és a keret kezdetének felismerését segíti; a hibadetektálásra az FCS szolgál, amely CRC-32 ellenőrzést tartalmaz.

Az IEEE 802.3 keretben a kétbájtos mező \textbf{Length}, vagyis az LLC PDU hosszát adja meg. Az Ethernet II keretben ugyanitt \textbf{Type} mező van, amely a felsőbb szintű protokollt azonosítja, például IPv4-et vagy ARP-ot. A gyakorlatban az Ethernet II keretformátum nagyon elterjedt.

\begin{figure}[htbp]
\centering
\resizebox{\textwidth}{!}{%
\begin{tikzpicture}[font=\scriptsize]
  \def\h{0.8}
  \draw[fill=lightaccent] (0,0) rectangle (1.1,\h);      \node[align=center] at (0.55,0.4) {PA\\7 B};
  \draw[fill=warnbg] (1.1,0) rectangle (1.8,\h);          \node[align=center] at (1.45,0.4) {SFD\\1 B};
  \draw[fill=lightaccent] (1.8,0) rectangle (2.8,\h);     \node[align=center] at (2.3,0.4) {DA\\6 B};
  \draw[fill=lightaccent] (2.8,0) rectangle (3.8,\h);     \node[align=center] at (3.3,0.4) {SA\\6 B};
  \draw[fill=warnbg] (3.8,0) rectangle (4.8,\h);          \node[align=center] at (4.3,0.4) {Len vagy Type\\2 B};
  \draw[fill=black!5] (4.8,0) rectangle (7.6,\h);         \node[align=center] at (6.2,0.4) {Data\\46--1500 B};
  \draw[fill=black!10] (7.6,0) rectangle (8.4,\h);        \node[align=center] at (8.0,0.4) {Pad\\ha kell};
  \draw[fill=lightaccent] (8.4,0) rectangle (9.3,\h);     \node[align=center] at (8.85,0.4) {FCS\\4 B};
  \draw[decorate, decoration={brace, mirror, amplitude=4pt}] (1.8,-0.15) -- (9.3,-0.15) node[midway, below=4pt]{64--1518 bájt: DA-tól FCS-ig};
  \node[align=center] at (4.65,1.25) {IEEE 802.3: Length + LLC PDU\qquad Ethernet II: Type + Data};
\end{tikzpicture}%
}
\caption{IEEE 802.3/Ethernet keret fő mezői}
\label{fig:zv23_ethernet_keret}
\end{figure}


\begin{center}

\begin{tabularx}{\textwidth}{L{28mm}L{22mm}X}
\toprule
\textbf{Mező} & \textbf{Méret} & \textbf{Feladat} \\
\midrule
Preamble & 7 bájt & Szinkronizáció, váltakozó 1 és 0 bitek. \\
SFD & 1 bájt & Start of Frame Delimiter, a keret kezdetét jelzi. \\
DA & 6 bájt & Destination Address, cél MAC-cím. \\
SA & 6 bájt & Source Address, forrás MAC-cím. \\
Length / Type & 2 bájt & IEEE 802.3-nál hosszmező, Ethernet II-nél felsőbb protokolltípus. \\
Data + Pad & 46--1500 bájt & Hasznos adat; rövid adat esetén kitöltés a minimális keretméret miatt. \\
FCS & 4 bájt & Frame Check Sequence, CRC-32 alapú hibadetektálás. \\
\bottomrule
\end{tabularx}

\end{center}

\subsection{Minimális és maximális keretméret}

A klasszikus Ethernetben a minimális kerethossz a célcímtől az FCS végéig \textbf{64 bájt}. Ebbe nem számít bele a 7 bájtos preambulum és az 1 bájtos SFD. Az adatmező minimális hasznos mérete 46 bájt, ezért ha a felsőbb réteg ennél rövidebb adatot ad át, a keretet \textbf{pad}, vagyis kitöltő bájtok egészítik ki.

A hagyományos maximális Ethernet keretméret az FCS-sel együtt \textbf{1518 bájt}, amely 1500 bájtos adatmezőt jelent. VLAN címkézés esetén a keret ennél 4 bájttal hosszabb lehet. A minimális méret történeti oka a CSMA/CD ütközésérzékelés, a maximális méret pedig a késleltetés és a közeg túl hosszú lefoglalásának korlátozása.

\subsection{MAC-címek}

A klasszikusan használt Ethernet MAC-cím \textbf{48 bites}, azaz 6 bájtos cím. Általában hexadecimális bájtpárokkal írjuk, például \texttt{00:1A:2B:3C:4D:5E}. A MAC-cím az adatkapcsolati rétegben helyi hálózati címzésre szolgál.

A célcím lehet:

\begin{itemize}
  \item \textbf{unicast}: egyetlen állomásnak szól;
  \item \textbf{multicast}: állomások egy csoportjának szól;
  \item \textbf{broadcast}: minden állomásnak szól; Ethernetben a csupa egyes cím, vagyis \texttt{FF:FF:FF:FF:FF:FF}.
\end{itemize}

A cím első oktettjében külön bitek jelzik, hogy egyedi vagy csoportcímről, illetve globálisan kiosztott vagy lokálisan adminisztrált címről van-e szó. A globálisan kiosztott címeknél az OUI -- Organizationally Unique Identifier -- azonosítja a gyártói címtartományt.

\begin{center}

\begin{tabularx}{\textwidth}{L{25mm}L{42mm}X}
\toprule
\textbf{Címtípus} & \textbf{Példa / jelzés} & \textbf{Jelentés} \\
\midrule
Unicast & I/G bit = 0 & Egyetlen állomásnak szóló egyedi cím. \\
Multicast & I/G bit = 1 & Állomások egy csoportjának szóló cím. \\
Broadcast & \texttt{\scriptsize FF:FF:FF:FF:FF:FF} & Minden állomásnak szóló speciális csoportcím. \\
Globális cím & U/L bit = 0 & Gyártóhoz kötött, IEEE által kiosztott OUI-t tartalmazó cím. \\
Lokális cím & U/L bit = 1 & Helyileg adminisztrált cím, amelyet a hálózat üzemeltetője vagy szoftver állíthat be. \\
\bottomrule
\end{tabularx}

\end{center}

\section{Duplex Ethernet és kapcsolt Ethernet}

\subsection{Half duplex működés}

\textbf{Half duplex} Ethernetnél egy állomás egy adott időben vagy ad, vagy vesz, és a közös ütközési tartomány miatt CSMA/CD-re van szükség. Ilyen volt a klasszikus koaxiális Ethernet és a hubbal összekötött 10BASE-T vagy 100BASE-TX hálózat. A hub fizikai szinten ismétli a jeleket, ezért az összes rákapcsolt állomás ugyanahhoz az ütközési tartományhoz tartozik.

Half duplex környezetben a hálózat maximális kiterjedését nemcsak a kábel fizikai tulajdonságai, hanem az ütközésérzékeléshez szükséges oda-vissza terjedési idő is korlátozza.

\subsection{Full duplex működés}

\textbf{Full duplex} Ethernetnél két végpont között külön adási és vételi irány áll rendelkezésre, tipikusan kapcsoló és végállomás között. Ilyenkor nincs közös ütközési tartomány, ezért CSMA/CD-re nincs szükség. A végpontok egyszerre tudnak adni és venni.

Full duplex működésnél a kapcsolat hosszkorlátját már nem a CSMA/CD slot time, hanem a fizikai közeg és az adott PHY specifikáció határozza meg. A túlterhelés kezelésére az IEEE 802.3x \textbf{Pause Frame} típusú áramlásszabályozást definiálhat: a fogadó fél ideiglenesen megállíthatja az adót.

\subsection{Kapcsolt Ethernet}

A modern Ethernet tipikusan \textbf{kapcsolt Ethernet}. A switch adatkapcsolati rétegbeli eszköz, amely MAC-címtáblát épít, és a kereteket csak a megfelelő port felé továbbítja. Ezzel az ütközési tartományok portonként elkülönülnek. Full duplex linkeknél gyakorlatilag megszűnik a klasszikus CSMA/CD jelentősége, bár a történeti és szabványbeli megértéshez továbbra is fontos.

\begin{center}

\begin{tabularx}{\textwidth}{L{30mm}L{43mm}X}
\toprule
\textbf{Szempont} & \textbf{Half duplex Ethernet} & \textbf{Full duplex Ethernet} \\
\midrule
Adás és vétel & Egy időben vagy adás, vagy vétel. & Egyidejű adás és vétel lehetséges. \\
Ütközés & Lehet ütközés közös ütközési tartományban. & Nincs klasszikus ütközés pont--pont linken. \\
MAC módszer & CSMA/CD szükséges. & CSMA/CD nem szükséges. \\
Tipikus eszköz & Koax, hub, megosztott közeg. & Switch és végállomás közötti pont--pont link. \\
Hosszkorlát & A fizikai közeg mellett a slot time is korlátoz. & Főleg a fizikai közeg és PHY specifikáció korlátoz. \\
\bottomrule
\end{tabularx}

\end{center}

\section{IEEE 802.11 WLAN általános felépítése}

\subsection{WLAN és rádiós közeg sajátosságai}

Az \textbf{IEEE 802.11} a vezeték nélküli helyi hálózatok, vagyis WLAN-ok szabványcsaládja. A rádiós közeg miatt több lényeges eltérés van a vezetékes Ethernethez képest:

\begin{itemize}
  \item a jel terjedése kevésbé kiszámítható, visszaverődés, csillapítás és interferencia léphet fel;
  \item nem biztos, hogy minden állomás hallja az összes többi állomást;
  \item az állomás adás közben tipikusan nem tud megbízhatóan ütközést érzékelni;
  \item a hibaarány nagyobb lehet, ezért fontos a nyugtázás és a keretdarabolás lehetősége;
  \item mobilitás és energiatakarékosság is megjelenik követelményként.
\end{itemize}

Ezért a WLAN nem CSMA/CD-t, hanem \textbf{CSMA/CA-t} használ: nem az ütközés gyors felismerése a fő cél, hanem az ütközés valószínűségének csökkentése.

\subsection{BSS, AP, DS és ESS}

A 802.11 terminológiában a hálózat cellákra épül.

\begin{itemize}
  \item \textbf{BSS -- Basic Service Set}: egy alapvető cella, vagyis egymással kommunikálni képes állomások halmaza.
  \item \textbf{AP -- Access Point}: hozzáférési pont, amely az infrastruktúrás BSS központi eleme.
  \item \textbf{DS -- Distribution System}: az AP-kat és cellákat összekötő elosztórendszer, gyakran vezetékes Ethernet hálózat.
  \item \textbf{ESS -- Extended Service Set}: több BSS összekapcsolt rendszere, amelyben a mobil állomás cellák között vándorolhat.
\end{itemize}


\begin{figure}[htbp]
\centering
\resizebox{0.92\textwidth}{!}{%
\begin{tikzpicture}[font=\scriptsize, node distance=9mm]
  \tikzset{ap/.style={draw, rounded corners, fill=warnbg, align=center, minimum width=16mm, minimum height=7mm},
           sta/.style={draw, circle, fill=lightaccent, align=center, minimum size=8mm}}
  \draw[fill=black!3, draw=accent!60, rounded corners] (-2.1,-1.7) rectangle (2.1,1.7);
  \draw[fill=black!3, draw=accent!60, rounded corners] (4.0,-1.7) rectangle (8.2,1.7);
  \node at (0,1.45) {BSS 1};
  \node at (6.1,1.45) {BSS 2};

  \node[ap] (ap1) at (0,0) {AP1};
  \node[ap] (ap2) at (6.1,0) {AP2};
  \node[sta] (s1) at (-1.2,0.8) {STA};
  \node[sta] (s2) at (1.1,-0.9) {STA};
  \node[sta] (s3) at (4.9,-0.8) {STA};
  \node[sta] (s4) at (7.2,0.8) {STA};

  \draw[dashed] (s1) -- (ap1);
  \draw[dashed] (s2) -- (ap1);
  \draw[dashed] (s3) -- (ap2);
  \draw[dashed] (s4) -- (ap2);

  \draw[very thick, accent] (ap1) -- node[above]{DS -- Distribution System} (ap2);
  \draw[decorate, decoration={brace, amplitude=5pt}] (-2.1,-2.05) -- (8.2,-2.05) node[midway, below=5pt]{ESS: több BSS összekapcsolva, roaming lehetőséggel};
\end{tikzpicture}%
}
\caption{IEEE 802.11 infrastruktúrás hálózat: BSS, AP, DS és ESS}
\label{fig:zv23_wlan_bss_ess}
\end{figure}


\subsection{Az Access Point főbb funkciói}

Az \textbf{Access Point} a WLAN infrastruktúrás működésének kulcseleme. Főbb feladatai:

\begin{itemize}
  \item beacon keretek küldése, amelyek azonosítják a hálózatot és szinkronizációs információt adnak;
  \item új állomások beléptetésének támogatása: passive scanning, active scanning, authentication, association;
  \item keretek továbbítása a vezeték nélküli állomások és a DS között;
  \item cellán belüli állomások közötti forgalom továbbítása, ha szükséges;
  \item roaming támogatása ESS esetén;
  \item energiatakarékos állomások számára keretek ideiglenes tárolása;
  \item PCF esetén Point Coordinator szerep ellátása.
\end{itemize}

\begin{center}

\begin{tabularx}{\textwidth}{L{26mm}L{32mm}X}
\toprule
\textbf{Fogalom} & \textbf{Teljes név} & \textbf{Szerep} \\
\midrule
STA & Station & Vezeték nélküli állomás, például kliens, laptop vagy telefon. \\
BSS & Basic Service Set & Egy WLAN cella, állomások és infrastruktúrás módban egy AP közös halmaza. \\
AP & Access Point & Hozzáférési pont; cellát szervez, beaconöket küld, beléptetést és továbbítást végez. \\
DS & Distribution System & Az AP-kat összekötő elosztórendszer, gyakran vezetékes Ethernet. \\
ESS & Extended Service Set & Több BSS-ből álló nagyobb WLAN, roaming lehetőséggel. \\
Beacon & -- & Periodikus menedzsment keret, amely hálózati és szinkronizációs információkat hordoz. \\
Association & -- & A kliens logikai kapcsolódása egy AP-hoz. \\
\bottomrule
\end{tabularx}

\end{center}

\section{IEEE 802.11 közeghozzáférés: CSMA/CA és DCF}

\subsection{Miért nem CSMA/CD?}

A vezetékes Ethernetnél az ütközés érzékelhető a közegen mérhető jelalak vagy feszültségszint alapján. Rádiós hálózatban ez sokkal nehezebb, mert az állomás saját adása nagyságrendekkel erősebb lehet, mint a másik állomástól érkező jel. Szemléletesen: a rádiós állomás adás közben nem hall jól.

További gond a \textbf{rejtett állomás} probléma. Előfordulhat, hogy két állomás nem hallja egymást, de mindketten ugyanannak az AP-nak vagy vevőnek küldenének. Ekkor mindkét állomás szabadnak érzékelheti a csatornát, miközben a vevőnél ütközés történik.

\subsection{DCF: Distributed Coordination Function}

A 802.11 alapértelmezett közeghozzáférési módja a \textbf{DCF -- Distributed Coordination Function}. Ez elosztott, versengéses működés: nincs szükség központi ütemezőre ahhoz, hogy az állomások hozzáférjenek a csatornához.

DCF esetén a küldés leegyszerűsített menete:

\begin{enumerate}
  \item az állomás adás előtt figyeli a csatornát;
  \item ha a csatorna DIFS ideig tétlen, indulhat a küldés vagy a backoff visszaszámlálás;
  \item ha a csatorna foglalt, az állomás megvárja a felszabadulást, majd véletlen backoff értéket választ;
  \item a backoff számláló csak tétlen közeg alatt csökken, foglalt közegnél megáll;
  \item amikor a backoff nullára csökken, az állomás ad;
  \item a hibátlanul vevő állomás SIFS után ACK kerettel nyugtáz;
  \item ha nincs ACK, a küldő hibát vagy ütközést feltételez, növeli a contention window-t és újrapróbálkozik.
\end{enumerate}

A \textbf{SIFS} rövidebb, mint a DIFS, ezért az ACK, CTS vagy töredékfolytatás elsőbbséget kap az új versengő adásokkal szemben. A \textbf{contention window} növelése exponenciális visszatartási logikát valósít meg.


\begin{figure}[htbp]
\centering
\resizebox{0.95\textwidth}{!}{%
\begin{tikzpicture}[font=\scriptsize]
  \draw[thick] (0,0) -- (12,0) node[right]{idő};
  \draw[fill=black!10] (0.4,0.15) rectangle (2.1,0.65) node[midway]{korábbi adás};
  \draw[<->] (2.1,0.85) -- (3.2,0.85) node[midway, above]{DIFS};
  \draw[fill=lightaccent] (3.2,0.15) rectangle (5.4,0.65) node[midway]{backoff rések};
  \draw[fill=warnbg] (5.4,0.15) rectangle (8.2,0.65) node[midway]{DATA};
  \draw[<->] (8.2,0.85) -- (8.8,0.85) node[midway, above]{SIFS};
  \draw[fill=lightaccent] (8.8,0.15) rectangle (9.8,0.65) node[midway]{ACK};
  \node[draw, rounded corners, fill=black!5, align=center] at (6,-1.05) {A backoff csak tétlen közeg alatt csökken.\\ACK hiánya esetén új, nagyobb ablakból választott backoff következik.};
\end{tikzpicture}%
}
\caption{DCF időzítés: DIFS, véletlen backoff, DATA, SIFS és ACK}
\label{fig:zv23_dcf_idozites}
\end{figure}


\subsection{Valóságos és virtuális vivőfigyelés}

A DCF kétféle csatornafoglaltság-érzékelést használhat:

\begin{itemize}
  \item \textbf{valóságos vivőfigyelés}: az állomás fizikailag figyeli a rádiós közeget;
  \item \textbf{virtuális vivőfigyelés}: a keretek Duration mezői alapján az állomás NAV értéket tart nyilván.
\end{itemize}

A \textbf{NAV -- Network Allocation Vector} azt jelzi, hogy az állomás szerint mennyi idő van még hátra a csatorna felszabadulásáig. Ha egy állomás RTS, CTS, adat vagy ACK keretből időtartamot olvas ki, frissítheti a NAV-ját, és addig foglaltnak tekinti a közeget, akkor is, ha fizikailag éppen nem érzékel jelet.

\subsection{RTS/CTS és rejtett állomás probléma}

A \textbf{RTS/CTS} mechanizmus a rejtett állomás probléma enyhítésére szolgál. A küldő először rövid \textbf{RTS -- Request To Send} keretet küld a vevőnek, amelyben megadja a tervezett adatátvitel időtartamát. A vevő \textbf{CTS -- Clear To Send} kerettel válaszol, amely szintén tartalmazza a szükséges időtartamot.

Azok az állomások, amelyek az RTS-t vagy a CTS-t hallják, beállítják a NAV-jukat, és nem kezdenek adásba a megadott idő alatt. Így akkor is védhető az adatkeret, ha a küldőt nem mindenki hallja, de a vevő CTS-e több állomáshoz eljut. Az RTS/CTS többletforgalmat okoz, ezért rövid kereteknél vagy kis terhelésnél gyakran kikapcsolható vagy küszöbértékhez kötött.


\begin{figure}[htbp]
\centering
\resizebox{0.9\textwidth}{!}{%
\begin{tikzpicture}[font=\scriptsize, node distance=18mm]
  \tikzset{sta/.style={draw, circle, fill=lightaccent, minimum size=9mm, align=center},
           ap/.style={draw, rounded corners, fill=warnbg, minimum width=13mm, minimum height=8mm, align=center}}
  \node[sta] (a) {A};
  \node[ap, right=of a] (b) {Vevő\\AP};
  \node[sta, right=of b] (c) {C};
  \draw[thick] (a) -- node[above]{hallják egymást} (b);
  \draw[thick] (b) -- node[above]{hallják egymást} (c);
  \draw[dashed, red] (a) .. controls +(0,-1.2) and +(0,-1.2) .. node[below]{A és C nem hallják egymást} (c);

  \node[draw, rounded corners, fill=black!5, align=left, below=12mm of b, text width=72mm] (steps)
    {1. A RTS-t küld a vevőnek.\\2. A vevő CTS-t küld, amit C is hall.\\3. A és C a Duration mezők alapján NAV-ot állít.\\4. Az adatkeret védettebb, mert C nem kezd adásba.};
  \draw[arrow] (b) -- (steps);
\end{tikzpicture}%
}
\caption{Rejtett állomás probléma és RTS/CTS alapötlet}
\label{fig:zv23_rts_cts}
\end{figure}


\section{DCF és PCF üzemmódok}

\subsection{PCF: Point Coordination Function}

A \textbf{PCF -- Point Coordination Function} központosított, ütközésmentes hozzáférési mód. A Point Coordinator szerepét rendszerint az AP látja el. A PCF célja, hogy versengésmentes időszakokat hozzon létre, például időkorlátos forgalom számára.

A PCF működése röviden:

\begin{itemize}
  \item az AP/Point Coordinator időnként magasabb prioritással megszerzi a csatornát;
  \item beacon keretben vagy vezérlőinformációval jelzi a versengésmentes periódust;
  \item a többi állomás NAV-ja alapján nem verseng a közegért;
  \item az AP polling keretekkel sorban megszólítja a kiválasztott állomásokat;
  \item a lekérdezett állomás küldhet adatot, majd a PCF periódus végén az AP lezárja a versengésmentes szakaszt.
\end{itemize}

A PCF a DCF fölé épülő kiegészítő lehetőség. A gyakorlatban kevésbé elterjedt, mint a DCF, de vizsgán ismerni kell az elvét: a DCF elosztott és versengéses, a PCF központosított és versengésmentes.

\begin{center}

\begin{tabularx}{\textwidth}{L{37mm}L{45mm}X}
\toprule
\textbf{Szempont} & \textbf{DCF} & \textbf{PCF} \\
\midrule
Teljes név & Distributed Coordination Function & Point Coordination Function \\
Hozzáférés jellege & Elosztott, versengéses. & Központosított, versengésmentes periódusokat hoz létre. \\
Vezérlés & Az állomások CSMA/CA szerint versengenek. & A Point Coordinator, tipikusan az AP, pollinggal irányít. \\
Interframe space & DIFS után indul az általános versengés. & PIFS-t használhat, ezért megelőzheti a DIFS-t használó állomásokat. \\
Tipikus szerep & Alapértelmezett WLAN adatforgalom. & Időkorlátos forgalom elvi támogatása, kevésbé elterjedt. \\
Ütközéskezelés & ACK hiánya, backoff, CW növelése, RTS/CTS opció. & A CFP alatt a polling miatt nincs normál versengés. \\
\bottomrule
\end{tabularx}

\end{center}

\subsection{DCF és PCF összevetése}

DCF esetén minden állomás azonos alapelv szerint verseng a csatornáért. Ez egyszerű, robusztus és jól illeszkedik a löketszerű adatforgalomhoz. PCF esetén a központi koordinátor határozza meg, hogy a versengésmentes időszakban ki adhat. Ez jobban szabályozható, de bonyolultabb és AP-tól függ.

A 802.11 időzítési logikában a rövidebb interframe space nagyobb prioritást jelent. A SIFS a legfontosabb közvetlen válaszokhoz tartozik, például ACK és CTS. A PIFS-t a PCF Point Coordinator használhatja, ezért megelőzheti a DIFS-t használó DCF állomásokat. A DIFS után indulhat az általános versengés.

\section{Gyakori vizsgahibák és összefoglalás}

\subsection{Gyakori félreértések}

\begin{itemize}
  \item \textbf{LLC és MAC összekeverése.} Az LLC a felsőbb rétegek felé ad logikai adatkapcsolati szolgáltatást, a MAC a közeghez való hozzáférést szabályozza.
  \item \textbf{Statikus és determinisztikus azonosítása.} A statikus előre felosztja a csatornát; a determinisztikus dinamikus módszer is lehet, csak ütközésmentes hozzáférést ad.
  \item \textbf{CSMA/CD és CSMA/CA összekeverése.} Ethernetnél az ütközésérzékelés, WLAN-nál az ütközéselkerülés a lényeg.
  \item \textbf{MAC-cím és IP-cím összemosása.} A MAC-cím adatkapcsolati rétegbeli, helyi hálózati cím; az IP-cím hálózati rétegbeli logikai cím.
  \item \textbf{A preambulum beleszámítása a 64 bájtos minimális keretméretbe.} A 64 bájt a célcímtől az FCS végéig értendő; a preambulum és SFD nincs benne.
  \item \textbf{Full duplex Ethernetnél CSMA/CD feltételezése.} Full duplex pont--pont linken nincs ütközés, ezért CSMA/CD-re nincs szükség.
  \item \textbf{AP szerepének leszűkítése egyszerű rádiós ismétlőre.} Az AP beaconöket küld, beléptetést kezel, forgalmat továbbít, DS-kapcsolatot biztosít, és energiatakarékossági/roaming funkciókat is támogathat.
  \item \textbf{DCF és PCF szerepének felcserélése.} A DCF az elosztott CSMA/CA alapmód, a PCF központi, AP által vezérelt versengésmentes mód.
\end{itemize}

\subsection{Rövid vizsgaválasz-minta}

\begin{summarybox}
Az adatkapcsolati réteg IEEE 802 szemléletben LLC és MAC alrétegre bontható. A MAC alréteg feladata közös üzenetszórásos csatorna megosztása több állomás között. A csatornamegosztás lehet statikus, például FDM vagy TDM, illetve dinamikus. Dinamikus esetben beszélhetünk versengő, ütközéses módszerekről, mint ALOHA, CSMA, CSMA/CD és CSMA/CA, valamint determinisztikus, ütközésmentes módszerekről, mint a vezérjeles vagy lekérdezéses megoldások. Az IEEE 802.3 Ethernet klasszikusan CSMA/CD-t használ: az állomás adás előtt figyeli a csatornát, adás közben ütközést érzékelhet, ütközéskor jam jelet küld, majd exponenciális visszatartás után újrapróbálkozik. Az Ethernet keret fő mezői a preambulum, SFD, cél MAC-cím, forrás MAC-cím, hossz vagy típus mező, adat, esetleges pad és FCS. A MAC-cím tipikusan 48 bites, lehet unicast, multicast vagy broadcast. A modern kapcsolt full duplex Ethernetben nincs közös ütközési tartomány, ezért CSMA/CD-re nincs szükség. Az IEEE 802.11 WLAN rádiós közegre épül, ahol ütközésérzékelés helyett CSMA/CA-t használnak. Infrastruktúrás hálózatban a BSS-t AP vezeti, több BSS-ből ESS állhat össze. A DCF elosztott, versengéses CSMA/CA alapmód, ACK-kal, backoffal, NAV-val és szükség esetén RTS/CTS-sel. A PCF ezzel szemben AP által irányított, központosított, versengésmentes hozzáférési mód.
\end{summarybox}
